\section{Conclusion}
\label{sec:conclusion}

We presented a two-pass 3D reconstruction system for automated line-of-sight audits at urban intersections. Our approach combines static scene reconstruction using Pi3 with dynamic object detection and tracking using YOLOv8x and ByteTrack. The system processes synchronized multi-camera video to produce a unified 3D scene representation with both static geometry and tracked dynamic objects represented as canonical primitives.

\textbf{Key findings:}
\begin{itemize}
    \item Pi3 produces dense, geometrically consistent point clouds from uncalibrated multi-view images with sub-2-pixel reprojection error
    \item SegFormer semantic segmentation enables robust ground plane estimation for accurate 2D-to-3D projection
    \item ByteTrack maintains stable object identities with low fragmentation ($<$5\%) through occlusions
    \item Multi-camera fusion improves track completeness from 67\% to 95\% compared to single-camera tracking
    \item Sub-meter 3D localization accuracy is achieved through ground plane projection with multi-camera fusion
\end{itemize}

\textbf{Limitations and future work:}
The current system operates offline and requires approximately 6 minutes to process 5 minutes of video. Real-time operation would require optimized inference pipelines and potentially edge deployment. Additionally, our evaluation is limited to a single intersection; broader validation across diverse urban environments is needed.

Future directions include: (1) implementing the full ray-casting visibility analysis pipeline to quantify line-of-sight metrics, (2) extending to real-time operation for live safety monitoring, (3) incorporating temporal consistency constraints for smoother 3D trajectories, and (4) developing actionable safety recommendations based on identified visibility hazards.

Our work demonstrates that combining state-of-the-art multi-view reconstruction with modern detection and tracking methods provides a viable foundation for automated intersection safety analysis, with potential applications in urban planning, traffic engineering, and autonomous vehicle validation.
